\documentclass[11pt,a4paper]{article}
\usepackage[utf8]{inputenc}
\usepackage[T1]{fontenc}
\usepackage{lmodern}
\usepackage{geometry}
\usepackage{hyperref}
\usepackage{amsfonts, amsmath, amssymb, amsthm}
\usepackage{mathtools}
\usepackage{enumitem}
\usepackage{microtype}
\usepackage{longtable}

\setlist{leftmargin=2em,itemsep=0.25ex,topsep=0.5ex}
\allowdisplaybreaks
\hypersetup{hidelinks}
\geometry{left=1in, right=1in, top=1in, bottom=1in}

\theoremstyle{plain}
\newtheorem{theorem}{Theorem}[section]
\newtheorem{conjecture}[theorem]{Conjecture}
\newtheorem{corollary}[theorem]{Corollary}
\newtheorem{lemma}[theorem]{Lemma}
\newtheorem{proposition}[theorem]{Proposition}

\theoremstyle{definition}
\newtheorem{definition}[theorem]{Definition}
\newtheorem{remark}[theorem]{Remark}
\newtheorem*{problem*}{Open question}

\newcommand{\C}{\mathbb{C}}
\newcommand{\e}{\mathrm{e}}
\newcommand{\Q}{\mathbb{Q}}
\newcommand{\R}{\mathbb{R}}
\newcommand{\Z}{\mathbb{Z}}
\newcommand{\Alg}{\overline{\mathbb{Q}}}
\newcommand{\Logs}{\mathcal{L}}
\newcommand{\aLogs}{\widetilde{\mathcal{L}}}
\DeclareMathOperator{\trdeg}{trdeg}

\title{The Diaz modulus conjecture: a machine-checked spine\\
\large and the shape of its obstruction}
\author{Carlo Perassi}
\date{September 2026\\[2pt]\small Version 1.1, 23 September 2026}

\begin{document}
\maketitle

\begin{abstract}
Diaz asked in 2004 whether a non-zero complex number can be a logarithm of an algebraic
number and have algebraic modulus. We give a short formal development of the question. All but
one of the numbered statements below have been checked by machine, and the exception is flagged
where it appears; Appendix~A records where each proof lives, so that any claim here can be
verified rather than taken.

The results divide in two. The first part is a spine: the strong four exponentials conjecture
settles the question outright, by a single instantiation; a statement it implies, which we
call~$(S)$, settles the part of it that the elementary theory reduces to; and the value that a
counterexample to $(S)$ would exhibit cannot be a root of unity. The second part is negative, and is the
reason for the paper. We show that one natural line of attack --- assembling the certified data
of a hypothetical counterexample into a matrix of logarithms with vanishing determinant ---
provably fails on one branch, that a second branch is circular in a precise sense, and that the
remaining branch terminates in a theorem of 1973--74. When this note was first written that
theorem had never been formalised; it has been since, in the form the branch needs, and the case
it governs is proved with it. The natural interpolation route is blocked as well: no first-order
operator with algebraic coefficients has algebraic values there, and the interpolation determinant
does not see the candidate's arithmetic at all. What is left is two statements, both open
mathematics rather than formalisation, and we say exactly which. A third, which the development had
carried until now, turns out not to be needed, and it and the statement it pairs with cannot both
be false. Along the way the transcendence theorems now available formally yield a few
unconditional results about classical constants, among them that at least one of
$\e^{\pi^{2}}$ and $\e^{i\pi^{3}}$ is transcendental.
\end{abstract}

\section{Setting}

Write $\Alg$ for the algebraic numbers,
\[
  \Logs \;=\; \{\lambda\in\C : \e^{\lambda}\in\Alg\}
\]
for the logarithms of algebraic numbers, and $\aLogs$ for the $\Alg$-vector subspace of $\C$
spanned by $1$ together with $\Logs$. Complex conjugation is written $\bar{\phantom{u}}$, and
$\overline{\e^{u}}=\e^{\bar u}$, so $\Logs$ is stable under it.

\begin{conjecture}[Diaz, 2004~\cite{diaz-2004}]\label{conj:diaz}
For every $u\in\C$ with $u\neq0$ and $\lvert u\rvert\in\Alg$, the number $\e^{u}$ is
transcendental.
\end{conjecture}

We call $u$ a \emph{candidate} when it violates this: $u\neq0$, $\lvert u\rvert\in\Alg$, and
$\e^{u}\in\Alg$. A candidate carries exactly four pieces of certified data,
\[
  u\in\Logs,\qquad \bar u\in\Logs,\qquad 2\pi i\in\Logs,\qquad u\bar u=\lvert u\rvert^{2}\in\Alg ,
\]
and the last of these is a \emph{quadratic} relation between logarithms. That is the whole
source of the difficulty: the linear theory of logarithms --- Baker's theorem and its relatives
--- says nothing about products, and the theory that does speak about products of two
logarithms is the four exponentials circle. Section~\ref{sec:ceiling} makes this precise, and
everything after it is an attempt to ask for less.

Two classical inputs are used without proof. \emph{Hermite--Lindemann}: for algebraic
$a\neq0$, $\e^{a}$ is transcendental. \emph{Transcendence of $\pi$}, and hence of $\pi^{2}$.
Three further classical theorems enter in the results added in version 1.1, and all three are
machine-checked: the Gel'fond--Schneider theorem, in the formalisation of Karatarakis and
Wiedijk~\cite{karatarakis-wiedijk}; the six exponentials theorem; and the four exponentials
theorem in transcendence degree one (Section~\ref{sec:map}).

\begin{conjecture}[Strong four exponentials]\label{conj:sfe}
Let $x_{1},x_{2}$ be linearly independent over $\Alg$ and let $y_{1},y_{2}$ be linearly
independent over $\Alg$. Then at least one of the four products $x_{i}y_{j}$ lies outside
$\aLogs$.
\end{conjecture}

This is Waldschmidt's strengthening of the four exponentials conjecture, in which linear
independence is taken over $\Alg$ rather than over $\Q$. Both are open.

\section{The ceiling}\label{sec:ceiling}

\begin{theorem}\label{thm:sfe-implies-diaz}
Conjecture~\ref{conj:sfe} implies Conjecture~\ref{conj:diaz}.
\end{theorem}

\begin{proof}
Let $u\neq0$ with $\lvert u\rvert\in\Alg$, and suppose $\e^{u}\in\Alg$.

If $u$ is algebraic, Hermite--Lindemann gives that $\e^{u}$ is transcendental, a
contradiction. So $u$ is transcendental, and so is $\bar u$, since the conjugate of an
algebraic number is algebraic. Consequently $(1,u)$ and $(1,\bar u)$ are each linearly
independent over $\Alg$: a relation $a+bu=0$ with $a,b\in\Alg$ not both zero would force
either $b\neq0$ and $u=-a/b\in\Alg$, or $b=0$ and $a=0$.

Apply Conjecture~\ref{conj:sfe} to $x=(1,u)$ and $y=(1,\bar u)$. The four products are
\[
  1,\qquad \bar u,\qquad u,\qquad u\bar u=\lvert u\rvert^{2},
\]
and all four lie in $\aLogs$: the first by definition of $\aLogs$; the second and third because
$\e^{\bar u}=\overline{\e^{u}}$ and $\e^{u}$ are algebraic, so $u,\bar u\in\Logs$; and the
fourth because $\lvert u\rvert^{2}$ is algebraic and $\Alg\cdot1\subseteq\aLogs$. This
contradicts Conjecture~\ref{conj:sfe}.
\end{proof}

\begin{remark}\label{rem:schanuel}
Schanuel's conjecture also suffices, by a different route. If $u$ and $\bar u$ are linearly
independent over $\Q$, Schanuel at $n=2$ gives
$\trdeg_{\Q}\Q(u,\bar u,\e^{u},\e^{\bar u})\geq2$, hence $\trdeg_{\Q}\Q(u,\bar u)\geq2$ once
the two exponentials are algebraic; but $u\bar u\in\Alg$ is a non-trivial algebraic relation, so
that degree is at most $1$. If instead $u=q\bar u$ with $q\in\Q$, then $\lvert q\rvert=1$, so
$q=\pm1$ and $u^{2}=\pm\lvert u\rvert^{2}\in\Alg$, making $u$ algebraic; Hermite--Lindemann
finishes. Schanuel at $n=1$ is itself Hermite--Lindemann, so no separate input is needed.
\end{remark}

Theorem~\ref{thm:sfe-implies-diaz} fixes the ceiling of the whole subject. Any decomposition of
Conjecture~\ref{conj:diaz} into cases produces statements that are consequences of
Conjecture~\ref{conj:sfe}, and the useful question is not whether a case follows from it --- all
of them do --- but whether a case follows from something \emph{weaker}. The next section
exhibits one such statement.


\section{Quantisation and the period lattice}\label{sec:quant}

This section is self-contained: apart from the transcendence of $\pi$ and of $\pi^{2}$ it uses
no arithmetic input at all. It ends by collapsing an entire branch of the problem to a single
sentence about a single real variable.

Throughout, $u$ is off both coordinate axes, meaning $\operatorname{Re}u\neq0$ and
$\operatorname{Im}u\neq0$.

\begin{theorem}[Quantisation]\label{thm:quant}
Let $u$ be off both axes with $\e^{u}$ real. Then $\operatorname{Im}u\in\pi\Z\setminus\{0\}$,
and $\lvert u\rvert^{2}>\pi^{2}$.
\end{theorem}

\begin{proof}
Write $u=t+i\theta$. Then $\e^{u}=\e^{t}(\cos\theta+i\sin\theta)$, so $\e^{u}$ real forces
$\e^{t}\sin\theta=0$, and $\e^{t}\neq0$ gives $\sin\theta=0$, that is $\theta=k\pi$ for some
$k\in\Z$. Since $\operatorname{Im}u\neq0$ we have $k\neq0$. Then
$\lvert u\rvert^{2}=t^{2}+k^{2}\pi^{2}\geq t^{2}+\pi^{2}>\pi^{2}$, using $t\neq0$.
\end{proof}

The lower bound is the only effective statement in this development, and it sharpens: if the
ratio $\e^{u}/\overline{\e^{u}}$ has finite order $m$, the same argument applied to $u/m$ gives
$\lvert u\rvert^{2}>\pi^{2}/m^{2}$. What that sharpening is worth is settled by the next
statement, which says the whole family carries exactly one bit of information.

\begin{proposition}\label{prop:quant-orbit}
For $u$ off the axes, the quantisation conclusion holds across the entire orbit
$\{qu: q\in\Q^{\times}\}$ if and only if $\operatorname{Re}u\neq0$.
\end{proposition}

No algebraicity is used. The point is negative: a predicate that looks like a family of
increasingly precise constraints is equivalent to the hypothesis one already has, so no split of
the problem along $m$ can be honest.

\begin{theorem}[Fibres]\label{thm:fibre}
Let $\alpha\in\Alg^{\times}$. At most two of the logarithms of $\alpha$ have algebraic modulus.
\end{theorem}

\begin{proof}
The logarithms of $\alpha$ are $u+2\pi i n$, $n\in\Z$, for any one of them $u$. Writing
$\theta=\operatorname{Im}u$,
\[
  \lvert u+2\pi i n\rvert^{2}
  \;=\;\lvert u\rvert^{2}+4\pi n\theta+4\pi^{2}n^{2}.
\]
Suppose $n_{1},n_{2},n_{3}$ are distinct and all three give an algebraic modulus, hence an
algebraic squared modulus. Subtracting the $n_{2}$ equation from the $n_{1}$ one and dividing by
$4(n_{1}-n_{2})\neq0$ gives $\pi\theta+\pi^{2}(n_{1}+n_{2})\in\Alg$, and likewise
$\pi\theta+\pi^{2}(n_{1}+n_{3})\in\Alg$. Subtracting these,
$\pi^{2}(n_{2}-n_{3})\in\Alg$ with $n_{2}\neq n_{3}$, so $\pi^{2}\in\Alg$ --- false.
\end{proof}

\begin{corollary}\label{cor:translate}
If two distinct logarithms of $\alpha$ have algebraic modulus, say $u$ and $u+2\pi i n$ with
$n\neq0$, then $\pi\bigl(\operatorname{Im}u+n\pi\bigr)\in\Alg$.
\end{corollary}

\begin{proof}
Immediate from the displayed identity: the difference of the two squared moduli is
$4\pi n\theta+4\pi^{2}n^{2}$, algebraic, and dividing by $4n$ gives
$\pi\theta+\pi^{2}n=\pi(\theta+n\pi)\in\Alg$.
\end{proof}

Corollary~\ref{cor:translate} is worth reading twice. The condition it produces was introduced
elsewhere in this subject as a technical case distinction; what it says is that the fibre over
$\e^{u}$ contains a \emph{second} point of algebraic modulus. Theorem~\ref{thm:fibre} bounds how
often that can happen, and eliminating $\operatorname{Im}u$ between the two displayed relations
returns a quartic relation in $\pi$ --- which is the reason the class it cuts out has
transcendence degree one over $\Q$. The period lattice contributes one bit and no more.

\subsection*{One sentence}

\begin{theorem}\label{thm:one-relation}
The following are equivalent.
\begin{enumerate}[label=\textup{(\roman*)}]
  \item For every $u\neq0$ off both axes with $\lvert u\rvert\in\Alg$ and $\e^{u}$ real,
        $\e^{u}$ is transcendental.
  \item For every real $t\neq0$ with $\e^{t}\in\Alg$, the number $t^{2}+\pi^{2}$ is
        transcendental.
\end{enumerate}
\end{theorem}

\begin{proof}
Assume (ii) and let $u$ be as in (i) with $\e^{u}\in\Alg$. By Theorem~\ref{thm:quant},
$u=t+ik\pi$ with $t=\operatorname{Re}u\neq0$ and $k\in\Z\setminus\{0\}$. Put $v=u/k$. Then
$\e^{v}$ is a $k$-th root of $\e^{u}$, hence algebraic; $\operatorname{Im}v=\pi$;
$\operatorname{Re}v=t/k\neq0$; and
$\lvert v\rvert^{2}=(t/k)^{2}+\pi^{2}$, which is algebraic because $\lvert u\rvert^{2}$ is and
$k$ is rational. Also $\e^{v}$ is real up to sign, so $\e^{t/k}=\pm\e^{v}$ is algebraic.
Applying (ii) at $t/k$ makes $(t/k)^{2}+\pi^{2}$ transcendental, contradicting the previous
sentence.

Conversely assume (i) and let $t\neq0$ be real with $\e^{t}\in\Alg$ and suppose
$t^{2}+\pi^{2}\in\Alg$. Set $u=t+i\pi$. Then $\e^{u}=-\e^{t}$ is real and algebraic,
$\lvert u\rvert^{2}=t^{2}+\pi^{2}$ is algebraic, and $u$ is off both axes. This contradicts~(i).
\end{proof}

The $k$ in that proof carries no content, which is the substance of the theorem: an entire branch
of the problem, quantified over two real parameters and an integer, is one assertion about one
real variable. Its smallest instance is a question one can state without any of the preceding
apparatus.

\begin{problem*}
Is $\sqrt{(\log 2)^{2}+\pi^{2}}$ algebraic?
\end{problem*}

Nothing in this development answers it. Section~\ref{sec:ceiling} shows the strong four
exponentials conjecture does; Section~\ref{sec:boundary} explains why the available
unconditional theory does not reach it.

\subsection*{A second point in the fibre, and the statement \texorpdfstring{$(S)$}{(S)}}

Corollary~\ref{cor:translate} says what it takes for the fibre over $\e^{u}$ to contain a second
point of algebraic modulus. On the branch where the angle is irrational, $\operatorname{Im}u\notin\pi\Q$,
that condition does more than constrain $u$: it manufactures a new logarithm.

\begin{proposition}\label{prop:period-aligned}
Let $u\in\Logs$ with $\operatorname{Im}u\notin\pi\Q$, and suppose
$\pi\bigl(\operatorname{Im}u+r\pi\bigr)\in\Alg$ for some rational $r\neq0$. Then there is a
non-zero algebraic $\gamma$ with $\e^{\gamma/(i\pi)}$ algebraic.
\end{proposition}

\begin{proof}
$\Logs$ is a $\Q$-vector space --- $\e^{(a/b)\lambda}$ is a root of $X^{b}-(\e^{\lambda})^{a}$ ---
and it contains $u$, $\bar u$ and $2\pi i$. So it contains
$\lambda=\tfrac12(u-\bar u)+\tfrac r2(2\pi i)=i\bigl(\operatorname{Im}u+r\pi\bigr)$. Put
$\gamma=i\pi\lambda=-\pi\bigl(\operatorname{Im}u+r\pi\bigr)$. It is algebraic by hypothesis, and
non-zero because $\operatorname{Im}u=-r\pi$ is excluded; and $\gamma/(i\pi)=\lambda\in\Logs$.
\end{proof}

So on this branch a candidate would exhibit a non-zero algebraic multiple of $1/(i\pi)$ inside
$\Logs$. That is the statement the branch reduces to.

\begin{conjecture}[$(S)$]\label{conj:S}
For every non-zero algebraic $\gamma$, the number $\e^{\gamma/(i\pi)}$ is transcendental.
Equivalently, no non-zero algebraic multiple of $1/(i\pi)$ lies in $\Logs$.
\end{conjecture}

\begin{theorem}\label{thm:S-of-sfe}
Conjecture~\ref{conj:sfe} implies Conjecture~\ref{conj:S}.
\end{theorem}

\begin{proof}
Suppose $\gamma\neq0$ is algebraic and $\lambda=\gamma/(i\pi)\in\Logs$. Apply
Conjecture~\ref{conj:sfe} to $x=(1,\lambda)$ and $y=(1,i\pi)$. The four products are
$1$, $i\pi$, $\lambda$ and $\gamma$, all in $\aLogs$. Both pairs are $\Alg$-linearly independent,
since $i\pi$ is transcendental and so is $\lambda$, whose algebraicity would make
$i\pi=\gamma/\lambda$ algebraic. This contradicts Conjecture~\ref{conj:sfe}.
\end{proof}

Whether $(S)$ is strictly weaker than Conjecture~\ref{conj:sfe} is not known; what is known is that
it asks for much less, one number against a whole configuration. Its hypothetical counterexamples
are also more constrained than the value $\e^{\gamma/(i\pi)}$ suggests.

\begin{proposition}\label{prop:root-of-unity}
For every non-zero algebraic $\gamma$, the number $\e^{\gamma/(i\pi)}$ is not a root of unity.
\end{proposition}

\begin{proof}
If $\e^{\gamma/(i\pi)}=\e^{2\pi ik/n}$, then $\gamma/(i\pi)=2\pi i\,(k/n+m)$ for some $m\in\Z$, so
$\gamma=-2\pi^{2}(k/n+m)$. As $\gamma\neq0$, the rational $k/n+m$ is non-zero and $\pi^{2}$ is
algebraic --- false.
\end{proof}

The statement also splits, and not by a case distinction. The set
$S_{0}=\{\gamma\in\Alg:\gamma/(i\pi)\in\Logs\}$ is a $\Q$-vector space, and it is stable under
conjugation because $\overline{i\pi}=-i\pi$ gives $\bar\gamma/(i\pi)=-\overline{\gamma/(i\pi)}$. So
with $\gamma$ it contains $\operatorname{Re}\gamma$ and $i\operatorname{Im}\gamma$, and $(S)$ is
equivalent to its two axis halves: for real $\gamma\neq0$, the number $\e^{-i\gamma/\pi}$, which
lies on the unit circle and is not a root of unity, is transcendental; and for $\gamma=i\beta$ with
$\beta$ real, $\e^{\beta/\pi}$ is transcendental --- that is, $\pi\neq\beta/\log\alpha$ for real
algebraic $\beta\neq0$ and algebraic $\alpha>0$, $\alpha\neq1$. The reduction is machine-checked,
and both halves are open. They cannot both fail.

\begin{theorem}\label{thm:S0}
\begin{enumerate}[label=\textup{(\alph*)}]
  \item If $\gamma_{1},\gamma_{2}\in S_{0}$ and $\gamma_{2}\neq0$, then $\gamma_{1}/\gamma_{2}\in\Q$.
  \item Every element of $S_{0}$ is real or purely imaginary.
  \item At least one of the two halves of $(S)$ holds.
\end{enumerate}
\end{theorem}

\begin{proof}
(a) The numbers $\lambda_{j}=\gamma_{j}/(i\pi)$ lie in $\Logs$, $\lambda_{2}\neq0$, and their ratio
$\gamma_{1}/\gamma_{2}$ is algebraic. If it were irrational, Gel'fond--Schneider would make
$\e^{\lambda_{1}}=\e^{(\gamma_{1}/\gamma_{2})\lambda_{2}}$ transcendental. (b) By the conjugation
stability of $S_{0}$ and (a), $\bar\gamma=q\gamma$ with $q\in\Q$; comparing real and imaginary
parts, either $\operatorname{Re}\gamma=0$, or $q=1$ and $\operatorname{Im}\gamma=0$. (c) A non-zero
real element and a non-zero purely imaginary one would be rationally proportional by (a).
\end{proof}

So $S_{0}$ is either $\{0\}$ or a single rational line on one axis, and $(S)$ can fail, if at all,
in one direction only. The imaginary half has a concrete reading: its exceptions are the
$\beta=\pi\log\alpha$ that happen to be algebraic, and by (a) all such $\log\alpha$ are rationally
proportional. For instance, at least one of $\pi\log2$ and $\pi\log3$ is transcendental, although
neither is known to be.

Which direction matters for Conjecture~\ref{conj:diaz} is settled by looking at what the branch
produces. In Proposition~\ref{prop:period-aligned} the number $\gamma=-\pi(\operatorname{Im}u+r\pi)$
is real. The formal decomposition reaches $(S)$ along one other branch, where $\pi\operatorname{Im}u$
itself is algebraic, and there $\gamma=-\pi\operatorname{Im}u$ is real as well. Both branches have
been reduced, by machine-checked reductions, to the real half alone. The imaginary half is
therefore not needed anywhere in the development (Section~\ref{sec:what}).



\section{The precise open boundary}\label{sec:boundary}

Attach to a candidate $u$, with $\rho=u\bar u$ and $r=\sqrt{\rho}$, the matrix
\[
  H(u,r)\;=\;\begin{pmatrix} u & r\\ r & \bar u\end{pmatrix},
  \qquad \det H(u,r)=u\bar u-r^{2}=0 .
\]
Its rows are dependent over $\C$. The next statement says they are dependent in no weaker
sense, and it is the sharpest way to see where this problem sits.

\begin{theorem}[No vanishing coefficient]\label{thm:novan}
Let $K\subseteq\C$ be a subfield with $r\in K$, $r\neq0$, and let $u$ be transcendental over
$K$ with $u\bar u=r^{2}$. Then for all non-zero $w,v\in K^{2}$,
\[
  w^{\mathsf T}H(u,r)\,v \;\neq\; 0 .
\]
\end{theorem}

\begin{proof}
Expanding, $w^{\mathsf T}Hv=w_{0}(uv_{0}+rv_{1})+w_{1}(rv_{0}+\bar u v_{1})$. Using
$\bar u=r^{2}/u$ and clearing the denominator $u\neq0$, the vanishing of this expression is a
quadratic relation for $u$ over $K$ with coefficients built from $w,v,r$; its leading coefficient
is $w_{0}v_{0}$ and its constant term is $r^{2}w_{1}v_{1}$. A vanishing of the expression makes
$u$ a root of a polynomial over $K$, so by transcendence that polynomial is identically zero:
$w_{0}v_{0}=0$, $r^{2}w_{1}v_{1}=0$, and the middle coefficient $r(w_{0}v_{1}+w_{1}v_{0})=0$.
With $r\neq0$ the first two give $w_{0}v_{0}=w_{1}v_{1}=0$, and combining with the third forces
$w=0$ or $v=0$.
\end{proof}

At $K=\Alg$ the transcendence hypothesis is exactly what a candidate supplies: an algebraic $u$
with $\e^{u}$ algebraic contradicts Hermite--Lindemann, so a candidate is transcendental.

So $H(u,r)$ is singular over $\C$ while every coefficient extracted from it over $\Alg$ is
non-zero. A candidate hides behind exactly this discrepancy: rank drops over $\C$ and nothing
detects it over $\Alg$.

One might hope this configuration is an artefact of a particular choice of matrix. It is not.

\begin{theorem}[Uniqueness of the obstruction]\label{thm:pencil}
Let $K\subseteq\C$ be a subfield, $u$ transcendental over $K$ with $u\neq0$ and $u\bar u\in K$,
and let $A,B,C$ be $2\times2$ matrices over $K$ with $\det(A+uB+\bar uC)=0$. Then there is
$c\in K$ with
\[
  \det(xA+yB+zC)\;=\;c\,\bigl(yz-(u\bar u)\,x^{2}\bigr)
  \qquad\text{for all }x,y,z\in\C .
\]
\end{theorem}

Every singular pencil over the candidate's hull is therefore the same object: its determinant is
a $K$-multiple of the single quadratic form $yz-\rho x^{2}$. Concretely, $H(u,r)$ factors as
\[
  \begin{pmatrix}0&r\\ r&0\end{pmatrix}
  \bigl(I+uB_{0}+\bar uC_{0}\bigr),\qquad
  B_{0}^{2}=C_{0}^{2}=0,\qquad \operatorname{tr}(B_{0}C_{0})=\rho^{-1},
\]
so the whole family is a nilpotent pencil pinned by one trace. There is no second configuration
to look for, and no gain in enlarging the search.

\subsection*{Descent to a conic, and what is missing}

Theorem~\ref{thm:pencil} sends the problem to the affine conic
\[
  X_{\rho}\;:\;XY=\rho ,
\]
on which a candidate is the point $(u,\bar u)\in\Logs^{2}$. This makes the boundary explicit,
and it also makes a tempting route collapse.

\begin{proposition}\label{prop:roy}
Fix $\rho\neq0$ and let $L\subseteq\C$ be any set. Suppose every point of $X_{\rho}\cap L^{2}$
lies in a $\Q$-rational vector subspace of $\C^{2}$ contained in $X_{\rho}$. Then
$X_{\rho}\cap L^{2}=\varnothing$.
\end{proposition}

\begin{proof}
A vector subspace contains the origin, and $0\cdot0=0\neq\rho$, so no vector subspace is
contained in $X_{\rho}$. The hypothesis can only be satisfied vacuously.
\end{proof}

The statement is trivial and the reason to record it is what it exposes when read backwards.
Roy's conjecture in the shape just used --- points of a variety with coordinates in $\Logs$ lie
on rational subspaces contained in it --- is available for \emph{homogeneous} quadrics, where
there are non-trivial subspaces to be placed on. The affine conic offers none, so the same
hypothesis becomes an emptiness statement rather than a structural one, and the gap between the
two is not a technicality but the location of the whole difficulty.

Stated in the form that is actually available: Proposition~12.25 of Waldschmidt's book, for
arbitrary $k\subseteq K$ and a $k$-subspace $\mathcal{E}$ of $K$, gives the conclusion for a
quadric defined over $k$ provided every non-zero matrix $M$ over $\mathcal{E}$ satisfies
\[
  \operatorname{rank}(M)\;>\;\tfrac12\,r_{\mathrm{str}}(M),
\]
the structural rank taken over $k$. Two things separate that from Diaz's conjecture, and they are
of different kinds.

\begin{enumerate}[label=\textup{(\alph*)}]
  \item \textbf{The rank inequality.} Over $(\Q,\C,\Logs)$ it is supplied at transcendence
        degree one by Theorem~15.30 of the same book --- which is why the
        transcendence-degree-one statements in this area are unconditional, and equally why they
        are not new. Over $(\Alg,\C,\aLogs)$, the coefficient field this problem requires, only
        the non-strict bound of Corollary~12.22 is known. This is being worked on; the
        Structural Rank Conjecture of Dasgupta and Kakde asserts something stronger.
  \item \textbf{Homogeneity.} Proposition~12.25 covers homogeneous quadrics; $X_{\rho}$ is
        affine. Waldschmidt makes the point himself, observing that the inhomogeneous extension
        would give the transcendence of $\e^{\lambda^{2}}$, and that $\e^{\pi^{2}}$ remains
        open. This is not being worked on.
\end{enumerate}

The distinction is worth keeping. An expectation that the subject will drift towards Diaz's
conjecture is reasonable for~(a) and unsupported for~(b), and it is~(b) that
Theorem~\ref{thm:one-relation} runs into: $\sqrt{(\log 2)^{2}+\pi^{2}}$ asks about a
\emph{sum}, and every available tool is homogeneous.

What the homogeneous theory does give at $\e^{\pi^{2}}$ is worth recording, because it is
unconditional and machine-checked. The four exponentials theorem in transcendence degree one
(Section~\ref{sec:map}) applies to the simplest singular matrix whose entries are tied by one ratio.

\begin{proposition}\label{prop:geometric}
Let $w\neq0$ and $z\notin\Q$ with $\trdeg_{\Q}\Q(w,z)\le1$. Then $w$, $wz$ and $wz^{2}$ are not
all in $\Logs$.
\end{proposition}

\begin{proof}
Otherwise the matrix $\begin{pmatrix}w&wz\\ wz&wz^{2}\end{pmatrix}$ has vanishing determinant and
entries in $\Logs$ generating transcendence degree at most one, so its rows or columns satisfy a
non-trivial rational relation $p\,w+q\,wz=0$; this forces $z=-p/q\in\Q$.
\end{proof}

\begin{corollary}\label{cor:epi2}
At least one of $\e^{\pi^{2}}$ and $\e^{i\pi^{3}}$ is transcendental.
\end{corollary}

\begin{proof}
Take $w=i\pi$ and $z=-i\pi$: then $\e^{w}=-1$, $wz=\pi^{2}$ and $wz^{2}=-i\pi^{3}$.
\end{proof}

This is as close to $\e^{\pi^{2}}$ as the homogeneous theory comes.



\section{Two barriers, and one branch that returns}\label{sec:map}

Section~\ref{sec:boundary} located the difficulty. This section records what is known about
methods rather than about the problem: two results saying that whole classes of approach cannot
work, and one observation about decomposition itself.

\subsection*{The candidate is invisible to algebraic coefficients}

\begin{theorem}[Indistinguishability]\label{thm:indist}
Let $u\neq0$ with $\e^{u}$ algebraic, and let $r\in\Alg$ be real and non-zero with
$u\bar u=r^{2}$. Then there exist a transcendental $t$ with $t\bar t=r^{2}$ and a ring
homomorphism $\Phi:\C\to\C$ fixing $\Alg$ pointwise, carrying $u$ to $t$, commuting with
conjugation on $\Alg(u)$, such that for every matrix $M$ over $\Alg(u)$ and all vectors $w,v$
over $\Alg$,
\[
  w^{\mathsf T}Mv=0 \iff w^{\mathsf T}\Phi(M)v=0 .
\]
\end{theorem}

The point $t$ is an ordinary point of the same circle --- $r\e^{i}$ will do --- with no
arithmetic distinction whatever. So no vanishing statement with algebraic coefficients,
formulated over the candidate's own hull, separates a hypothetical candidate from a point that
is certainly not one. Any argument that would refute a candidate by exhibiting such a vanishing
relation refutes nothing.

That the intertwining with conjugation comes for free is worth a line: since $\bar u=\rho/u$,
conjugation is a rational function of the generator on $\Alg(u)$, so any $\Phi$ matching
generators commutes with it automatically. The symmetry is not arranged; it is forced.

\subsection*{The threshold is a dimension}

\begin{theorem}[Dimension count]\label{thm:dim}
Let $u$ be a candidate. Then $1,u,\bar u$ all lie in $\aLogs$ and, granting
Hermite--Lindemann, they span an $\Alg$-subspace of dimension exactly $3$. There exist
$\Alg$-independent pairs $x=(x_{1},x_{2})$ and $y=(y_{1},y_{2})$ with all four products
$x_{i}y_{j}$ in that span; and for no $\Alg$-independent $x$ of length $2$ and $y$ of length $3$
do all six products $x_{i}y_{j}$ lie in it.
\end{theorem}

The certificate a candidate offers is three-dimensional. A $2\times2$ configuration needs four
products and fits; a $2\times3$ configuration needs six and cannot. That is the whole reason the
strong four exponentials conjecture reaches this problem while the six exponentials
\emph{theorem} --- which is proved, and which would otherwise be the natural instrument --- does
not. The obstruction is not depth but dimension, and no sharpening of six exponentials that
stays inside a three-dimensional certificate space can help. Where there is room the theorem
works as expected: with $x=(1,\pi)$ and $y=(\log2,\log3,\log5)$ it shows that at least one of
$2^{\pi}$, $3^{\pi}$, $5^{\pi}$ is transcendental.

\subsection*{A decomposition that returns to its start}

Splitting the conjecture into cases produces statements that are formally weaker, and it is
tempting to read a long chain of such splits as progress. It need not be.

\begin{proposition}\label{prop:loop}
Let $(\mathrm{NT})$ be the assertion: for every $u$ with $\e^{u}$ and $\e^{\bar u}$ algebraic,
$\operatorname{Re}u\neq0$, $\operatorname{Im}u\neq0$, and $\operatorname{Re}u$,
$\operatorname{Im}u$, $\operatorname{Re}u/\operatorname{Im}u$ all transcendental, the modulus
$\lvert u\rvert$ is transcendental. Then, granting Hermite--Lindemann, $(\mathrm{NT})$ is
equivalent to Conjecture~\ref{conj:diaz}.
\end{proposition}

\begin{proof}
Assume $(\mathrm{NT})$ and let $u\neq0$ have $\lvert u\rvert\in\Alg$ and $\e^{u}\in\Alg$. Then
$\e^{\bar u}=\overline{\e^{u}}$ is algebraic. If any of $\operatorname{Re}u$,
$\operatorname{Im}u$, $\operatorname{Re}u/\operatorname{Im}u$ is algebraic, or if either of
$\operatorname{Re}u$, $\operatorname{Im}u$ vanishes, then $\lvert u\rvert^{2}\in\Alg$ forces $u$
itself algebraic --- in each case one coordinate is algebraic and the identity
$\lvert u\rvert^{2}=(\operatorname{Re}u)^{2}+(\operatorname{Im}u)^{2}$ supplies the other ---
and Hermite--Lindemann applies. Otherwise $(\mathrm{NT})$ makes $\lvert u\rvert$ transcendental,
contradicting the hypothesis. The converse is immediate: a $u$ witnessing the failure of
$(\mathrm{NT})$ has algebraic modulus and algebraic $\e^{u}$.
\end{proof}

$(\mathrm{NT})$ is recorded formally and is open; the equivalence just proved is not itself
formalised, and Appendix~\ref{app:certs} marks it so.

$(\mathrm{NT})$ is a natural-looking restriction of the conjecture, reached by imposing
genericity on a conjugate pair, and it sits several case splits below the original statement. It
is nonetheless the original statement. A decomposition can descend a long way and arrive back
where it began, and only closing the cycle detects this; reading the chain forwards never does.

\subsection*{What is not available}

Two standard instruments are inapplicable here for reasons that are structural rather than
quantitative, and it is worth recording which.

\emph{Gel'fond--Schneider} produces transcendence of $\alpha^{\beta}$ from an algebraic
$\alpha\neq0,1$ and an algebraic irrational $\beta$. Applied to the configuration above it
manufactures transcendence at a point unrelated to the one in question; the number whose
transcendence is wanted is not of that shape. It is not idle elsewhere: it is what confines the
exceptions to $(S)$ to one rational line (Theorem~\ref{thm:S0}).

\emph{Baker's theorem} bounds non-vanishing linear forms in logarithms from below. The relevant
form here vanishes identically by hypothesis --- that is what being a candidate means --- so the
bound has nothing to act on. Using it would require knowing in advance that the form is
non-zero, which is the conclusion sought.

Finally, the transcendence-degree-one case of four exponentials, which does apply on part of the
locus, is a theorem of 1973--74 (Brownawell; Waldschmidt), recorded as Theorem~1 of
Roy--Waldschmidt~\cite{roy-waldschmidt-1995}. When this note was first written it was available
for citation and for nothing else, and the obstacle to a formal development looked like
Philippon's zero estimate together with the apparatus around it. That reading was wrong about
the route. Waldschmidt's own proof~\cite{waldschmidt-1973}, with the toolbox of his 1971
paper~\cite{waldschmidt-1971}, uses none of it: Siegel's lemma, a Gel'fond-type transcendence
criterion, the maximum principle, and a zero count for exponential polynomials that needs no
separation hypothesis --- the last being what removes Baker's theorem from the 1973 argument.
Along that route the theorem is now formalised, and the case of the branch that it governs is
proved with it; see Appendix~\ref{app:certs}. What remains on that branch is not the 1973--74
theorem but two statements about $1/\pi$, both open. The formal statement takes all four exponentials to be algebraic,
which is more than the 1973 theorem assumes and exactly what this setting supplies.

\subsection*{Interpolation on the lattice sees nothing}

The classical way to attack $\e^{u}$ and $\e^{\bar u}$ together is to run the two-variable
exponential system
\[
  F(z,w)=\exp(uz+\bar u\,w)
\]
on the lattice $\Z^{2}$, where its values $F(m,n)=\e^{mu}\e^{n\bar u}$ are algebraic, and to feed
an auxiliary function built from it to an interpolation or Schneider--Lang argument. Two
elementary facts stand in the way.

\begin{proposition}[No first-order arithmetic operator]\label{prop:no-deriv}
Let $u$ be a candidate and $X=a\,\partial_{z}+b\,\partial_{w}$ with $a,b$ polynomials with
algebraic coefficients. If $(XF)(m,n)$ is algebraic for every $(m,n)\in\Z^{2}$, then $a=b=0$.
\end{proposition}

\begin{proof}
$(XF)(m,n)=\bigl(a(m,n)u+b(m,n)\bar u\bigr)F(m,n)$, and $F(m,n)$ is a non-zero algebraic
number, so $a(m,n)u+b(m,n)\bar u$ is algebraic. A candidate lies on neither axis --- on an axis,
$|u|\in\Alg$ would make $u$ algebraic, against Hermite--Lindemann --- so $u$ and $\bar u$ are
$\Q$-linearly independent, as in the proof of Theorem~\ref{thm:line}, and Baker's theorem in its
inhomogeneous form gives $a(m,n)=b(m,n)=0$. A polynomial in two variables vanishing on $\Z^{2}$
is zero.
\end{proof}

The same holds for $a,b$ rational functions with algebraic coefficients regular on $\Z^{2}$, by
clearing denominators; the formal statement is the polynomial one. The consequence is that the
criterion of Schneider--Lang~\cite[Theorem~4.1]{waldschmidt-book}, which asks the derivatives of
the functions to stay in a ring generated over a number field, has no input here. No operator of
order one has algebraic values on the lattice. The operator that does, $\partial_{z}\partial_{w}$,
with $\partial_{z}\partial_{w}F=u\bar u\,F=\rho F$, is exactly where the candidate's relation
enters, and it is not a derivation.

\begin{proposition}[Kronecker factorisation]\label{prop:kronecker}
Let $u$ be a candidate and $\alpha=\e^{u}$. For $0\le a,b,m,n\le N$ the matrix
$M=\bigl(\exp(aum+b\bar u n)\bigr)$, with rows indexed by $(a,b)$ and columns by $(m,n)$, is the
Kronecker product $A\otimes B$ of $A=(\alpha^{am})$ and $B=(\bar\alpha^{\,bn})$. Hence
\[
  \det M=(\det A)^{N+1}(\det B)^{N+1}\neq0 .
\]
\end{proposition}

\begin{proof}
$\exp(aum+b\bar un)=\alpha^{am}\bar\alpha^{\,bn}=A_{a,m}B_{b,n}$, which is the Kronecker product,
and $\det(A\otimes B)=(\det A)^{N+1}(\det B)^{N+1}$. Each factor is a Vandermonde determinant, in
the nodes $\alpha^{a}$ and $\bar\alpha^{\,b}$ respectively, and the nodes are pairwise distinct
because $|\alpha|=\e^{\operatorname{Re}u}\neq1$, a candidate having $\operatorname{Re}u\neq0$.
\end{proof}

So the interpolation determinant of the natural system splits into two independent one-variable
determinants, and neither $\rho$ nor the relation $u\bar u=\rho$ appears in it. Schwarz' lemma for
Cartesian products~\cite[Proposition~4.7]{waldschmidt-book} asks for vanishing on a full simplex
of derivatives, which throws the method back on the values; and the values, by
Proposition~\ref{prop:kronecker}, see nothing of the candidate's arithmetic.

\section{Exclusions outside the hull}\label{sec:outside}

Section~\ref{sec:map} showed that nothing inside $W_{u}=\Alg+\Alg u+\Alg\bar u$ decides the
question: the hull a candidate certifies out of itself is three-dimensional, and the templates
that would refute it need four dimensions. The way past that is to leave the hull --- to adjoin
one further number and ask what it would have to be. This section collects what that yields.
Everything here is a necessary condition on a hypothetical counterexample; none of it is an
emptiness proof, and the final subsection says where the method stops.

Throughout, $u$ is a candidate, $\rho=|u|^{2}\in\Alg^{\times}$, and $\bar u=\rho/u$, so that
$u^{-1}=\bar u/\rho$ lies in $\aLogs$ along with $u$ and $1$.

\subsection*{Algebraic lines, and distance to an algebraic point}

\begin{theorem}\label{thm:line}
Let $\lambda\in\Logs$ with $\operatorname{Re}\lambda\neq0$ and $\operatorname{Im}\lambda\neq0$.
There are no $B\in\Alg^{\times}$ and $C\in\Alg$ with
\[
  B\lambda+\bar B\,\bar\lambda+C=0 .
\]
\end{theorem}

\begin{proof}
Off both axes, $\lambda$ and $\bar\lambda$ are $\Q$-linearly independent: taking real and
imaginary parts of $p\lambda+q\bar\lambda=0$ with $p,q\in\Q$ gives
$(p+q)\operatorname{Re}\lambda=0$ and $(p-q)\operatorname{Im}\lambda=0$, whence $p=q=0$. Both lie
in $\Logs$, since $\e^{\bar\lambda}=\overline{\e^{\lambda}}$ is algebraic. Baker's theorem in its
inhomogeneous form then makes the non-zero $\Alg$-combination $B\lambda+\bar B\bar\lambda$
transcendental, so it cannot equal the algebraic number $-C$.
\end{proof}

The condition has a geometric reading. Writing $B=B_{1}+iB_{2}$ and $\lambda=x+iy$, the relation
is $2(B_{1}x-B_{2}y)+C=0$: an affine line with real algebraic coefficients, not both zero. So
Theorem~\ref{thm:line} says a logarithm off the axes lies on no such line. For a candidate the
same conclusion is available with less input, and it is worth recording which:

\begin{remark}\label{rem:line-cheap}
For a candidate $u=x+iy$, no relation $Ax+By=C$ with $A,B,C\in\Alg\cap\R$ and $(A,B)\neq(0,0)$
holds, by Hermite--Lindemann alone. Such a line meets the circle $X^{2}+Y^{2}=\rho$ only in
points with algebraic coordinates --- solve the linear equation for one coordinate and substitute,
the resulting quadratic having non-zero leading coefficient because $A,B$ are real --- and $u$ is
transcendental. No linear forms in logarithms are needed. This form now has a machine-checked
proof of its own; Theorem~\ref{thm:line} remains the stronger statement, dropping the hypothesis
on $|u|$ at the cost of Baker's theorem. It is also the degree-one case of
Proposition~\ref{prop:kernel} below: a non-zero polynomial of degree at most one cannot be
divisible by $X^{2}+Y^{2}-\rho$.
\end{remark}

\begin{corollary}\label{cor:distance}
For every $a\in\Alg^{\times}$, $\ |u-a|\notin\Alg$. Equivalently, a candidate lies on no circle
of algebraic radius about a non-zero algebraic centre. This statement is now machine-checked.
\end{corollary}

\begin{proof}
If $|u-a|$ were algebraic then so would be
$2\operatorname{Re}(\bar a u)=|u|^{2}+|a|^{2}-|u-a|^{2}$, which is a forbidden line by
Remark~\ref{rem:line-cheap}. The exclusion of $a=0$ is not cosmetic: $|u|$ is algebraic by
hypothesis.
\end{proof}

\subsection*{The multiplier module}

The next statement determines, completely, which numbers can multiply a candidate back into
$\aLogs$. Its input is Roy's strong six exponentials theorem: for $\Alg$-linearly independent
$x_{1},x_{2}$ and $\Alg$-linearly independent $y_{1},y_{2},y_{3}$, at least one of the six
products $x_{i}y_{j}$ lies outside $\aLogs$.

\begin{theorem}\label{thm:multiplier}
Assume the strong six exponentials theorem. For every candidate $u$,
\[
  \{z\in\aLogs:\ uz\in\aLogs\}=\Alg+\Alg\,u^{-1}. \tag{M}
\]
\end{theorem}

\begin{proof}
For the inclusion $\supseteq$: $1$ and $u^{-1}$ lie in $\aLogs$, and for $z=a+bu^{-1}$ with
$a,b\in\Alg$ one has $uz=au+b\in\aLogs$.

For $\subseteq$, suppose $z\in\aLogs$, $uz\in\aLogs$ and $z\notin\Alg+\Alg u^{-1}$. Since $u$ is
transcendental, $1,u$ are $\Alg$-linearly independent, and so are $1,u^{-1}$; the assumption on
$z$ extends the second pair to an $\Alg$-linearly independent triple $z,1,u^{-1}$. The six
products of $(1,u)$ with $(z,1,u^{-1})$ are
\[
  z,\qquad 1,\qquad u^{-1},\qquad uz,\qquad u,\qquad u\cdot u^{-1}=1,
\]
all of them in $\aLogs$. This contradicts the strong six exponentials theorem.
\end{proof}

The force of (M) is that the right-hand side is closed under nothing useful. Two instances:

\begin{corollary}\label{cor:square}
Assume the strong six exponentials theorem. For every candidate $u$ and all $a\in\Alg^{\times}$,
$b,c\in\Alg$,
\[
  u^{2}\notin\aLogs, \qquad\text{hence}\qquad au^{2}+bu+c\notin\aLogs .
\]
In particular $\e^{au^{2}+bu+c}$ is transcendental.
\end{corollary}

\begin{proof}
Take $z=u$ in (M). If $u^{2}\in\aLogs$ then $u=a'+b'u^{-1}$ for some $a',b'\in\Alg$, so
$u^{2}-a'u-b'=0$ and $u$ is algebraic. The second form follows because $\aLogs$ contains $u$ and
$\Alg$ and is an $\Alg$-subspace.
\end{proof}

\begin{corollary}\label{cor:recip}
Assume the strong six exponentials theorem. For every candidate $u$ and every $a\in\Alg^{\times}$,
\[
  \frac{1}{u-a}\notin\aLogs .
\]
\end{corollary}

\begin{proof}
Put $z=(u-a)^{-1}$, a legitimate test element because $uz=1+az$, so $uz\in\aLogs$ as soon as
$z$ does. Then (M) gives $z=c+du^{-1}$ with $c,d\in\Alg$, and clearing denominators,
\[
  c\,u^{2}+(d-ca-1)\,u-da=0 .
\]
The coefficients are algebraic and not all zero: $c=0$ forces $d=1$, and then the constant term
is $-a\neq0$. So $u$ is algebraic, a contradiction.
\end{proof}

The contrast is the point. A candidate's reciprocal $1/u$ necessarily lies in $\aLogs$, while no
algebraically shifted reciprocal does.

\subsection*{Commensurable pairs, unconditionally}

Corollary~\ref{cor:square} rests on Roy's theorem, which has no formal proof. Part of it survives
on the four exponentials theorem in transcendence degree one alone, through a statement about
pairs of logarithms.

\begin{theorem}\label{thm:pair}
Let $u,v\in\Logs\setminus\{0\}$ with $|u|^{2}/|v|^{2}\in\Q$, and suppose
$\trdeg_{\Q}\Q(u,\bar u,v,\bar v)\le1$. Then $v\in\Q u$ or $v\in\Q\bar u$.
\end{theorem}

\begin{proof}
With $c=|u|^{2}/|v|^{2}$, the matrix $\begin{pmatrix}u&v\\ c\bar v&\bar u\end{pmatrix}$ has
determinant $|u|^{2}-c|v|^{2}=0$ and non-zero entries in $\Logs$. A column relation gives
$v\in\Q u$, a row relation $v\in\Q\bar u$.
\end{proof}

The hypotheses are satisfiable, by $u=\log2$ and $v=\log4$, and both alternatives occur. For a
candidate, $\bar u=\rho/u$, so the transcendence-degree hypothesis holds for every $v$ algebraic
over $\Q(u)$.

\begin{corollary}\label{cor:monomial}
Let $u$ be a candidate, $k\ge2$, and $c\in\Alg^{\times}$ with $|c|^{2}\rho^{k-1}\in\Q$. Then
$\e^{cu^{k}}$ is transcendental; in particular $\e^{u^{k}/|u|^{k-1}}$ is, for every $k\ge2$.
\end{corollary}

\begin{proof}
If $v=cu^{k}\in\Logs$, then $|v|^{2}=(|c|^{2}\rho^{k-1})\rho$, and Theorem~\ref{thm:pair} gives
$cu^{k}=qu$ or $cu^{k}=q\rho/u$ with $q\in\Q$. Either way a positive power of $u$ is algebraic, so
$u$ is algebraic, against Hermite--Lindemann.
\end{proof}

Negative exponents follow by applying the corollary to $\bar u$, since
$cu^{-k}=(c/\rho^{k})\,\bar u^{k}$.

\subsection*{Saturation at one logarithm}

\begin{theorem}\label{thm:saturation}
Let $u$ be a candidate and $\ell\in\Logs$. If $u\in\Alg+\Alg\ell$ then $u\in\Q\ell$.
\end{theorem}

\begin{proof}
Write $u=a+b\ell$ with $a,b\in\Alg$. If $u$ and $\ell$ are $\Q$-linearly independent, Baker's
inhomogeneous theorem makes $u-b\ell$ transcendental, and it equals $a$. So they are
$\Q$-dependent: $pu+q\ell=0$ with $p,q\in\Q$ not both zero. Here $p\neq0$, since $p=0$ would give
$q\ell=0$ and hence $\ell=0$, making $u=a$ algebraic against Hermite--Lindemann. Thus
$u=(-q/p)\ell$.
\end{proof}

So inside an algebraic line spanned by one logarithm, the elements of $\Logs$ are exactly the
\emph{rational} multiples of that logarithm: the algebraic span collapses to the rational one. The
instance $\ell=i\pi$ closes off a plane one might have expected to have to examine.

\begin{corollary}\label{cor:api}
No candidate has the form $a+b\pi$ with $a,b\in\Alg$.
\end{corollary}

\begin{proof}
$i\pi\in\Logs$, and $a+b\pi=a+(-ib)(i\pi)$ with $-ib$ algebraic. By
Theorem~\ref{thm:saturation}, $u=r\,i\pi$ for some $r\in\Q$, necessarily non-zero; then
$|u|=|r|\pi$ is transcendental, contradicting the candidate's algebraic modulus.
\end{proof}

The boundary of this argument should be stated explicitly, because it is easy to over-read.
Theorem~\ref{thm:saturation} handles one logarithm at a time, with the candidate itself as the
second. It says nothing about a combination of two independent logarithms, and in particular it
does not touch $\log2+i\pi$, a logarithm of $-2$ and the smallest open test case of
\S\ref{sec:quant}. The distance between the two situations is the whole difficulty.

\subsection*{Generic on its own circle}

One further exclusion belongs with these, and it is the strongest of them. It was the last
statement in this note to acquire a machine-checked proof, and the route the formalisation took
is not the one given below.

\begin{proposition}\label{prop:kernel}
Let $u=x+iy$ be a candidate and $\rho=|u|^{2}$. For every $P\in\Alg[X,Y]$,
\[
  P(x,y)=0 \quad\Longleftrightarrow\quad X^{2}+Y^{2}-\rho \ \text{ divides } \ P .
\]
\end{proposition}

\begin{proof}
The implication $\Leftarrow$ is immediate from $x^{2}+y^{2}=\rho$. For $\Rightarrow$, make the
invertible linear change of variables $U=X+iY$, $V=X-iY$ over $\Alg$, under which $X^{2}+Y^{2}$
becomes $UV$; let $Q\in\Alg[U,V]$ be the image of $P$, so that $Q(u,\bar u)=P(x,y)=0$ and
$\bar u=\rho/u$. Modulo $UV-\rho$ the ring $\Alg[U,V]$ becomes $\Alg[U,U^{-1}]$, and the residue
of $Q$ is a Laurent polynomial vanishing at $u$. Since $u$ is transcendental, a non-zero Laurent
polynomial over $\Alg$ cannot vanish at $u$: clearing denominators would give an algebraic
relation for $u$. Hence the residue is zero, $Q\in(UV-\rho)$, and $P\in(X^{2}+Y^{2}-\rho)$.
\end{proof}

A candidate is therefore algebraically generic on its own circle: it lies on no algebraic plane
curve that does not contain that entire circle. It lies on no parabola, and on no ellipse or
hyperbola other than through the circle appearing as a component. What the proposition does not
do is eliminate the circle itself, which is precisely where the difficulty has been all along ---
so the statement is sharp in form and empty in consequence for the main question.

The formal proof takes a different route from the one above, and the difference is instructive.
Rather than passing to the coordinate ring of the conic, it reads $P$ in $\Alg[X][Y]$ and divides
by the monic quadratic $Y^{2}+(X^{2}-\rho)$, leaving a remainder $A(X)Y+B(X)$ which must vanish at
$(x,y)$. Two facts finish it: $x$ and $y$ are each transcendental --- if either were algebraic,
$x^{2}+y^{2}=\rho$ would make the other algebraic and hence $u$ algebraic --- and $y\notin\Alg(x)$,
because $\rho-X^{2}$ is squarefree of degree two and so not a square in $\Alg(X)$. The change of
variables is the shorter argument on paper; the division is the one a proof assistant prefers,
because it needs no quotient ring.

\section{What would have to be proved}\label{sec:what}

Three statements, in increasing strength, each of which settles Conjecture~\ref{conj:diaz}.

\begin{enumerate}[label=\textup{(\arabic*)}]
  \item The single relation of Theorem~\ref{thm:one-relation}, for the branch it governs: for
        real $t\neq0$ with $\e^{t}$ algebraic, $t^{2}+\pi^{2}$ is transcendental. Smallest open
        instance: is $\sqrt{(\log2)^{2}+\pi^{2}}$ algebraic?
  \item The rank inequality of \S\ref{sec:boundary}(a) over $(\Alg,\C,\aLogs)$, together with an
        inhomogeneous extension of Proposition~12.25. The first component is under active study;
        the second is not.
  \item The strong four exponentials conjecture, which settles the whole question at once by
        Theorem~\ref{thm:sfe-implies-diaz}.
\end{enumerate}

Nothing here brings any of the three closer. What the development does establish is where the
question is not: not in a better matrix, by Theorem~\ref{thm:pencil}; not in an algebraic
invariant of the candidate's hull, by Theorem~\ref{thm:indist}; not in the six exponentials
theorem, by Theorem~\ref{thm:dim}; and not, on one whole branch, in further case analysis, by
Proposition~\ref{prop:loop}.

\subsection*{Where the formal development stops}

The case analysis behind this note has been carried to its end. Every branch of
Conjecture~\ref{conj:diaz} is either closed by a machine-checked proof or reduced, by a
machine-checked reduction, to one of exactly two statements, and nothing else in it is open:
\begin{enumerate}[label=\textup{(\roman*)}]
  \item the real half of $(S)$: for real algebraic $\gamma\neq0$, $\e^{-i\gamma/\pi}$ is
        transcendental;
  \item the statement $(\mathrm{NT})$ of Proposition~\ref{prop:loop}, which is equivalent to the
        conjecture itself.
\end{enumerate}
Both follow from Conjecture~\ref{conj:sfe}. Neither is a formalisation problem. Version 1.0 of
this note listed a third, the imaginary half of $(S)$: for real algebraic $\beta\neq0$,
$\e^{\beta/\pi}$ is transcendental. It remains open, but the development no longer depends on it,
since both branches that reach $(S)$ produce a real $\gamma$. By Theorem~\ref{thm:S0}, it and the
real half cannot both fail.

It is worth saying why the four exponentials theorem in transcendence degree one, now proved, does
not reach them, since each lives in transcendence degree one: under a counterexample to $(S)$ the
numbers $\lambda$ and $i\pi$ are algebraically dependent, $\lambda=\gamma/(i\pi)$, and a
candidate's $u$ and $\bar u$ are algebraically dependent because $\bar u=\rho/u$. The obstruction is the input,
not the degree. The configurations that
Theorems~\ref{thm:sfe-implies-diaz} and~\ref{thm:S-of-sfe} use have products
$1,\ \bar u,\ u,\ \rho$ and $1,\ i\pi,\ \lambda,\ \gamma$, and in each two of the four are non-zero
algebraic numbers. A non-zero algebraic number is never in $\Logs$, by Hermite--Lindemann. The
four exponentials theorem, and the four exponentials conjecture above it, ask for their products
to be logarithms of algebraic numbers; only the strong form, whose products are allowed to range
over $\aLogs$, accepts these.
That distinction, $\Logs$ against $\aLogs$, is the whole distance between what is proved and what
is needed.

\begin{thebibliography}{10}
\bibitem{diaz-2004} G.~Diaz, \emph{Utilisation de la conjugaison complexe dans l'étude de la
transcendance de valeurs de la fonction exponentielle usuelle}, J. Théor. Nombres Bordeaux
\textbf{16} (2004), 535--553.
\bibitem{diaz-2007} G.~Diaz, \emph{Produits et quotients de combinaisons linéaires de logarithmes
de nombres algébriques : conjectures et résultats partiels}, J. Théor. Nombres Bordeaux
\textbf{19} (2007), 373--391. It combines Roy's strong six exponentials theorem with complex
conjugation to obtain cases of the strong four exponentials conjecture; Roy's theorem is not
formalised, so those results are not used here.
\bibitem{karatarakis-wiedijk} M.~Karatarakis and F.~Wiedijk, \emph{A formalization of the
Gelfond-Schneider theorem}, arXiv:2603.24823 (2026).
\bibitem{roy-waldschmidt-1995} D.~Roy and M.~Waldschmidt, \emph{Quadratic relations between
logarithms of algebraic numbers}, Proc. Japan Acad. Ser. A \textbf{71} (1995), 151--153.
\bibitem{brownawell-1974} W.~D.~Brownawell, \emph{The algebraic independence of certain numbers
related to the exponential function}, J. Number Theory \textbf{6} (1974), 22--31.
\bibitem{waldschmidt-1973} M.~Waldschmidt, \emph{Solution du huitième problème de Schneider},
J. Number Theory \textbf{5} (1973), 191--202.
\bibitem{waldschmidt-1971} M.~Waldschmidt, \emph{Indépendance algébrique des valeurs de la
fonction exponentielle}, Bull. Soc. Math. France \textbf{99} (1971), 285--304.
\bibitem{waldschmidt-book} M.~Waldschmidt, \emph{Diophantine Approximation on Linear Algebraic
Groups}, Grundlehren der mathematischen Wissenschaften \textbf{326}, Springer, 2000.
\bibitem{dasgupta-kakde} S.~Dasgupta and M.~Kakde, \emph{On ranks of matrices of logarithms},
preprint.
\end{thebibliography}



\appendix
\section{Machine-checked certificates}\label{app:certs}

Every numbered statement above but one has been checked in a formal development, in Lean~4
against a pinned revision of its mathematical library. The table records where each one lives, so
that any claim in this note can be verified rather than taken on trust. The environment is stated
here once and nowhere else.

The proofs exist in two places, and the difference matters to a reader who wants to check them.
Each identifier below names a node on a collaborative proving platform, where the statement, the
submitted proof and its verdict are public. Every proved statement is also mirrored in a
self-contained Lean project at
\begin{center}\nolinkurl{https://github.com/carlok/diaz-modulus-lean}\end{center}
which builds from source in one command, depends on no platform, and assumes nothing beyond Lean's
own three axioms. Where the two differ in naming, the project keeps the platform's spelling.

This is version 1.1 of the note; version 1.0 of 21 September 2026 is the first stable release.
Every row of the table below was checked against the platform by script on 23 September 2026, and
the mirrored project held all 182 proved results of the development at that date.

The 1973--74 theorem's row stands for a development rather than a single proof: the criterion of
Waldschmidt 1971, the zero count for exponential polynomials, the construction of 1973 with its
auxiliary function, extrapolation and norm step, and the assembly above them. They are separate
nodes under the \nolinkurl{FourExp} namespace on the platform, and separate modules in the
project.

Two further statements, not numbered above because the note does not use them directly, are
proved alongside Proposition~\ref{prop:kernel} and belong to its argument: the real and imaginary
parts of a candidate are each transcendental
(\nolinkurl{DiazModulus.candidate_re_transcendental},
\nolinkurl{DiazModulus.candidate_im_transcendental}).

A \textbf{Proved} row means the statement itself has a machine-checked proof. An \textbf{Open}
row means the statement is recorded formally but not proved; every such row here is an open
conjecture. A \textbf{Not formalised} row means what it says: the statement is proved in this note and has no formal counterpart at all.
Proposition~\ref{prop:loop} is flagged that way, and its row needs a word: the identifier given
is the statement $(\mathrm{NT})$, which is recorded formally and is Open; the \emph{equivalence}
proved in Proposition~\ref{prop:loop} is what has no formal counterpart.

\begin{center}
\footnotesize
\begin{longtable}{@{}llp{6.2cm}@{}}
\textbf{Result} & \textbf{Status} & \textbf{Identifier} \\
\hline
\endhead
Conj.~\ref{conj:diaz} & Open (conjecture) & \nolinkurl{DiazModulus.diaz_modulus_conjecture} \\
Thm~\ref{thm:sfe-implies-diaz} & Proved & \nolinkurl{DiazModulus.diaz_of_sfe} \\
Rem.~\ref{rem:schanuel} & Proved & \nolinkurl{DiazModulus.diaz_of_schanuel} \\
Thm~\ref{thm:quant} & Proved & \nolinkurl{Diaz.real_quantisation} \\
Prop.~\ref{prop:quant-orbit} & Proved & \nolinkurl{Diaz.quantisation_orbit_iff_re_ne_zero} \\
Thm~\ref{thm:fibre} & Proved & \nolinkurl{Diaz.fibre_at_most_two} \\
Cor.~\ref{cor:translate} & Proved & \nolinkurl{Diaz.q_translate_unique} \\
Thm~\ref{thm:one-relation} & Proved & \nolinkurl{DiazModulus.leaf_iff_one} \\
Thm~\ref{thm:novan} & Proved & \nolinkurl{Diaz.no_vanishing_coeff} \\
Thm~\ref{thm:pencil} & Proved & \nolinkurl{Diaz.det_pencil_eq_conic} \\
Prop.~\ref{prop:roy} & Proved & \nolinkurl{Diaz.roy_conic_implies_empty} \\
Thm~\ref{thm:indist} & Proved & \nolinkurl{Diaz.candidate_indistinguishable_by_coeff} \\
Thm~\ref{thm:dim} & Proved & \nolinkurl{DiazModulus.sixExponentials_cannot_refute_candidate} \\
Prop.~\ref{prop:loop} & Not formalised & \nolinkurl{DiazModulus.norm_transcendental_of_generic_conj_pair} (Open) \\
\S\ref{sec:map}, 1973--74 thm & Proved & \nolinkurl{DiazModulus.four_exponentials_trdeg_one} \\
Prop.~\ref{prop:period-aligned} & Proved & \nolinkurl{DiazModulus.recip_pi_log_of_period_aligned} \\
Conj.~\ref{conj:S} & Open (conjecture) & \nolinkurl{DiazModulus.recip_pi_not_log} \\
Thm~\ref{thm:S-of-sfe} & Proved & \nolinkurl{DiazModulus.recip_pi_not_log_of_sfe} \\
Prop.~\ref{prop:root-of-unity} & Proved & \nolinkurl{DiazModulus.recip_pi_exp_value_not_root_of_unity} \\
Thm~\ref{thm:S0}(a) & Proved & \nolinkurl{DiazModulus.recip_pi_log_rational_line} \\
Thm~\ref{thm:S0}(b) & Proved & \nolinkurl{DiazModulus.recip_pi_log_on_axis} \\
Thm~\ref{thm:S0}(c) & Proved & \nolinkurl{DiazModulus.recip_pi_not_log_real_or_imag} \\
\S\ref{sec:quant}, $\pi\log2$ or $\pi\log3$ & Proved & \nolinkurl{DiazModulus.pi_log_two_or_pi_log_three_transcendental} \\
\S\ref{sec:quant}, real half of $(S)$ & Open (conjecture) & \nolinkurl{DiazModulus.recip_pi_not_log_real_gamma} \\
\S\ref{sec:quant}, imaginary half of $(S)$ & Open (conjecture) & \nolinkurl{DiazModulus.recip_pi_not_log_imag_gamma} \\
Prop.~\ref{prop:no-deriv} & Proved & \nolinkurl{DiazModulus.no_first_order_arithmetic_operator} \\
Prop.~\ref{prop:kronecker} & Proved & \nolinkurl{DiazModulus.kronecker_factorisation} \\
Prop.~\ref{prop:geometric} & Proved & \nolinkurl{DiazModulus.geometric_triple_not_logs} \\
Cor.~\ref{cor:epi2} & Proved & \nolinkurl{DiazModulus.exp_pi_sq_or_exp_i_pi_cube_transcendental} \\
\S\ref{sec:map}, $2^{\pi}$, $3^{\pi}$ or $5^{\pi}$ & Proved & \nolinkurl{DiazModulus.two_three_five_pow_pi_transcendental} \\
Thm~\ref{thm:line} & Proved & \nolinkurl{DiazModulus.no_algebraic_generalized_line} \\
Rem.~\ref{rem:line-cheap} & Proved & \nolinkurl{DiazModulus.candidate_no_real_algebraic_line} \\
Cor.~\ref{cor:distance} & Proved & \nolinkurl{DiazModulus.candidate_distance_transcendental} \\
Thm~\ref{thm:multiplier} & Proved & \nolinkurl{DiazModulus.candidate_multiplier_module} \\
Cor.~\ref{cor:square} & Proved & \nolinkurl{DiazModulus.candidate_multiplier_module} \\
Cor.~\ref{cor:recip} & Proved & \nolinkurl{DiazModulus.candidate_multiplier_module} \\
Thm~\ref{thm:pair} & Proved & \nolinkurl{DiazModulus.log_pair_rigid_of_trdeg_one} \\
Cor.~\ref{cor:monomial} & Proved & \nolinkurl{DiazModulus.candidate_monomial_not_log} \\
Thm~\ref{thm:saturation} & Proved & \nolinkurl{DiazModulus.candidate_one_log_saturation} \\
Cor.~\ref{cor:api} & Proved & \nolinkurl{DiazModulus.candidate_one_log_saturation} \\
Prop.~\ref{prop:kernel} & Proved & \nolinkurl{DiazModulus.candidate_vanishing_ideal} \\
Hermite--Lindemann & Proved & \nolinkurl{DiazModulus.hermite_lindemann_holds} \\
\end{longtable}
\end{center}

Two entries need a word. Hermite--Lindemann is used freely in the text as classical; it is listed
because the library used does not contain it, so it had to be established rather than cited. The
same holds for Gel'fond--Schneider, which is proved on the platform in another mission, from the
formalisation of Karatarakis and Wiedijk, and is mirrored in the project. And
the statement $(\mathrm{NT})$ of Proposition~\ref{prop:loop} was contributed by another
participant in the same development, independently of the work reported here; the equivalence
with Conjecture~\ref{conj:diaz} is what is established above.

Some results in the development are not represented here at all. Roughly a third of it consists
of closure lemmas, coercions and bookkeeping with no content at this level of exposition, and a
further group restates in one set of coordinates what another states in a different one. Neither
belongs in a paper, and neither is listed.


\end{document}
